\chapter{\textit{SCRIBE}: Reliable Distributed Model Learning Under Limited Communication} \label{chap:scribe}

\section{Introduction}
In Figure~\ref{fig:high-level-process-diagram}, we identified three necessary components to the autonomous exploration problem. The first of the three components, online data-driven model learning, is an area of rich study. The field of machine learning, including both reinforcement and non-reinforcement learning, has been a powerful driving factor in better understanding data-driven modeling of unknown spaces; the current literature on online model learning centers on utilizing gathered observations in conjunction with machine learning techniques to learn models during exploration. In particular, autonomous exploration literature typically attempts to combine model learning methods with statements on model confidence; these uncertainty estimates both enables more robust validation of learning and provides guidance on where the autonomous explorer still lacks environment understanding. Examples of uncertainty-aware modeling techniques include Gaussian Processes and other Bayesian machine learning approaches \cite{viseras2016decentralized, allamraju2017communication}.

Distributing the model learning process, however, is a non-trivial process, particularly in ensuring all agents in the distributed setting have models that are consistent with each other. Current distributed implementations in practice either directly share gathered observations across agents, with each agent independently learning a model on the collective set of observations, or average model parameters across agents. Typical model learning techniques grow exponentially in computational complexity, making the former approach computationally taxing; observation sharing also struggles when observations are large and the communication bandwidth is limited. The latter set of approaches, in conjunction with learned model approximations \cite{jang2020multi}, reduces the computational load on individual agents, but provides no guarantees for consistency of model certainty. Parameter averaging typically overestimates model certainty when communication is intermittent \cite{chen2002estimation}, further reducing the effectiveness of this latter set of techniques.

In this chapter, we discuss the problem of distributing learning of environment models across multiple distributed agents under limiting communication conditions, while guaranteeing that this shared model is reliably accurate, consistent, and understood by all agents. We shall propose, test, and validate our algorithmic solution, termed SCRIBE: Signal-aware Collaboration for Reliable Information-based Estimation. Specifically, we combine a lightweight, linear formulation for environment model learning, with a linear model update process that is consistent in confidence while robust to intermittent networks. We utilize a linear fixed-dimensional approximation of Gaussian Processes, alongside a hybrid linear filter interleaving an Information Kalman Filter with distributed Covariance Intersection. Agents can repeatedly update their learned models with any other available agents without loss of consistency, removing the need for all agents to be simultaneously involved during model updates.

\subsection{Related Work}

Online distributed model learning underwater is a challenging problem, as the communication necessary for even the base decentralization is both complex and restrictive. The limitations on effective decentralization - the difficulty of propagating messages across an agent-agent communication network, the restrictions on computation, and the bandwidth and noise present in transmitted messages - has resulted in comparatively limited contemporary work on online distributed learning underwater. However, this learning problem has also been addressed across a variety of non-ocean related fields, among them decentralized motion tracking and filtering, distributed sensor placement and fusion, and distributed adaptive sampling. The same paradigm is broadly followed across application fields, involving acquiring observations, learning a model from the gathered data, and finally disseminating the learned knowledge across the multi-agent or multi-sensor network. Depending on the system being modeled, the intended uses of the learned models, and the fidelity of the model structure, the specifics of the approach taken can vary widely. Crucially, aspects of the challenges faced underwater are observable in various other application fields, proving a rich bed of potential partial solutions for the ocean domain.

In the broader literature, there have been several approaches to learning system models, from linear filter models to reinforcement learning techniques, with an overview further discussed in the Appendix. These approaches range from model-based learning, where agents know the model structure but are yet unsure of the parameters, to be inferred via data, to model-free learning, where the structure of the model is either unknown or otherwise not leveraged and a generic function approximator is instead trained on existing data. The former set of approaches corresponds to mechanisms like Bayesian Model Inference and Expectation-Minimization, while the latter set of approaches corresponds to mechanisms including Q-learning, Policy Gradients, and evolutionary algorithms. These approaches have traditionally been developed in single-agent systems and centralized systems, requiring observations be gathered, and learning to happen, at a single node.

Distributed variants of these methods have struggled with balancing computational feasibility, learning accuracy, and robustness to overconfidence. Decentralization has typically succeeded by prioritizing model parameter inference over model function generation. Multi-agent inference methods, which use estimators such as filters or probabilistic graphical models, require the model structure to be known by all agents \textit{a priori}. This narrows the focus from learning the underlying model function, to learning the estimator's parameter set; for simplicity, we can term this as \textit{distributed estimation}. Filter-based methods are typically computationally efficient but have hefty requirements on communication, thus struggling on limited and intermittent communication networks; other consensus-based methods, such as Metropolis-Hastings Markov Chain averaging, provide similar benefits along with similar pitfalls. Both classes of approaches also suffer from overconfident predictions, particularly under intermittent and noisy networks. In comparison, the other popular classes of model learning, probabilistic graphical models and Covariance Intersection, can infer parameters that produce consistent predictions. However, depending on the underlying estimator complexity, graphical models and Covariance Intersection approaches can further suffer from high computation and communication requirements.

Estimator-based inference has enjoyed widespread adoption because of its simplification of the distribution problem. By constraining the potential parameter space, agents need only infer over a subset of parameters, enabling use of inference techniques that assume some level of model structure. This reduces the need of peer-to-peer consensus to only agreeing on parameter values; comparatively, function generation (model-free learning) requires agents first agree on the model structure itself before further inferring and consenting on the model parameter values. However, the trade-off is that estimator-based inference has less flexibility in system representation, sometimes significantly so. Recent developments in distributing function generation has seen promise, by exploiting the static update step in Gaussian Processes (GPs). While GPs are generic function approximators, i.e. while GPs attempt to generically learn any functional relation between data via only data, the GP model is updated via a static functional process. \cite{allamraju2017communication} shows that it is possible to exploit this update structure, along with artificially generated observations, to perform distributed update and consensus for model learning. However, this approach lacks the consistency guarantee provided by estimator-based methods like probabilistic graphical models and Covariance Intersection.

We propose a new, hybrid modeling approach that blends the flexibility of Gaussian Processes with the consistency, efficiency, and robustness of estimator-based methods. To achieve this, we note existing developments in the field on estimator-based inference that is consistent and communication-robust, and on approximating Gaussian Processes as estimator-based inference. Accordingly, we reformulate the distributed model learning problem as decentralized updates to an approximation of a Gaussian Process via robust and efficient estimator-based inference. Empirically, we note that with a sufficiently high-dimensional approximation, we retain the flexibility of a Gaussian Process while still effectively constraining the problem inference space for efficient distributed inference.

\cite{jadaliha2012environmental} demonstrates the linear approximation of a Gaussian Process. Recognizing that GPs perform nonlinear function approximation through linear regression by masking the nonlinearity via covariance kernels, \cite{jadaliha2012environmental} shows that the process can be approximately realized as the weighted linear composition of (potentially non-linear) basis functions. The basis functions will be determined prior to online model learning, and with a sufficiently high number of basis functions with good distribution parameter values, the resultant model learning demonstrates comparable flexibility to a Gaussian Process. \cite{jadaliha2012environmental, cui2015mutual} both show that this approximation can be efficiently updated given observations via a Kalman Filter-based, thus proving the approach as an efficient data-based model learning technique. However, the approach is still susceptible to intermittent communication networks and overconfidence in the learned model.

Comparatively, \cite{tamjidi2021unifying} demonstrates a hybrid Kalman Filter-Covariance Intersection approach for efficient, consistent, and communication-robust system model inference. The additive update of the  Information Kalman Filter (whose details are described further in the Appendix) easily lends to distribution, but in unreliable communication networks can become prone to inconsistency and divergence in learned system models across multiple agents. By introducing an efficient channel filter (Metropolis-Hastings Markov Chain filter update) to the additive update step for an Information Kalman Filter, they eliminate the need for a single agent to wait for the entire multi-agent network before updating, reducing the chance of signal corruption or loss. By further introducing Covariance Intersection on agent priors, this hybrid approach eliminates the need of stability across the entire communication network, providing robustness to intermittent communication. However, \cite{tamjidi2021unifying} accomplishes this by exploiting the linear update structure of the Kalman Filter, which accordingly requires linear system estimator formulations, constricting the application scope of the technique.

We recognize that the Gaussian Process approximation derived in \cite{jadaliha2012environmental} results in a linear system estimator, and can thus benefit from the technique proposed in \cite{tamjidi2021unifying}. Hence, we propose SCRIBE as the combination of the two methods, by first approximating an environment model as a linear combination of Gaussian Scalar Fields, and then updating this model online from observations gathered by agents operating in a decentralized setting.

\section{Model Formulation} \label{sec:env-model-formula}

We will assume the environment as a hidden Markov model, where we can only approximate the true system dynamics after a noisy observation process. Furthermore, we will further assume that we can satisfactorily approximate a stationary environment process as a linear update system driven by stochastic changes \cite{jadaliha2012environmental}. Given these assumptions, we can begin defining the system.

Let $x$ be the location of the vehicle and $\Omega$ the sampling region of interest such that $x\in\Omega\subseteq\mathbb{R}^3$. Let us also define the process we are measuring (ex. $\text{CO}_2$ concentration, etc.) as a scalar field $\mathcal{F}$ on $\Omega$ that is also time-varying, i.e. $\mathcal{F}(x,k):=\Omega,\mathbb{R}^+\rightarrow\mathbb{R}$. Following \cite{jadaliha2012environmental}, can more explicitly define $\mathcal{F}$ as a linear summation of $n_{\phi}$ weighted Gaussian functions:

\begin{equation}
    \mathcal{F}(x,k)=\sum_{i=1}^{n_\phi}\psi_i(x)\phi_i(k)=\bm{\psi}^T(x)\bm{\phi}(k)
    \label{eq:scalar-field-def}
\end{equation}

$\bm{\psi}^T(x)$ is the finite set of Gaussian functions and $\bm{\phi}(k)$ is the corresponding weight vector, where each $\psi_i(x)$ is a Gaussian function with corresponding parameters $\tau_i$, $\mu_i$, and $\sigma_i$ characterizing the distribution:

\begin{equation}
    \psi_i(x) = \frac{1}{\tau_i}\exp\left(-\frac{||x-\mu_i||^2}{\sigma_i^2}\right)
    \label{eq:gaussian-function-characterization}
\end{equation}

Note that we write $\bm{\phi}(k)$ as varying with time and not space, specifically varying with a discrete time step $k$.

During implementation, we can select $n_{\phi}$ variations of these parameter values that we find sufficiently cover the field we are trying to monitor. For example, if we set $n_{\phi}$ to be 5, we may set $\tau=0.01$ and $\sigma=0.1$ for all five Gaussian functions and vary $\mu_i$ over 5 different values equally spaced between -1 to 1. Because the set of Gaussian functions $\bm{\psi}(x)$ used in the environment are decided \textit{a priori}, the accuracy of the model depends on the quality of the weight vector $\bm{\phi}(k)$. The goal of observations is also to best approximate the weight vector.

We represent the time-varying nature of the scalar field $\mathcal{F}$ as being stochastically driven by some zero-mean additive white Gaussian noise $\mathbf{w}(k)\sim\mathcal{N}(\mathbf{0},\mathbf{Q}(k))$. $\mathbf{Q}(k)$ is the covariance of this stochastic noise process. We assume a lack of external driving input to the environment field. Assuming environment stationarity can allow us to make the further simplification that there are no inherent dynamics to the environment system; however, this is not a necessary assumption, and we will proceed without considering this additional assumption. This lets us write the change of $\bm{\phi}(k)$ over discrete time as:

\begin{equation}
    \bm{\phi}(k+1) = \mathbf{A}\bm{\phi}(k) + \mathbf{w}(k)
    \label{eq:field-weights-linear-update}
\end{equation}

where $\mathbf{A}=\mathbf{I}_{n_{\phi}}$ in the special case of no inherent environment process dynamics.

Let us now define the model for environment observations; we start by considering $n_r$ agents in the field, each capable of taking up to $n_s$ samples at varying locations between each discrete timestep model update (constrained by their specific kinematic dynamics). Each agent $j$ will take samples at locations $\{x_j^{[1]},\ x_j^{[2]},\ \hdots,\ x_j^{[n_s]}\}$. The corresponding observations made at each location can be written as $\{z_j(x_j^{[1]},k),\ z_j(x_j^{[2]},k),\ \hdots,\ z_j(x_j^{[n_s]},k)\}$, where $z_j(\cdot, k)$ represents the observation process of agent $j$ at timestep $k$. We model observations as noisy observations of the true scalar process, giving us the following definition for an observation at a single location $x_j^{[l]}$:

\begin{align}
    z_j(x_j^{[l]},k) &= \mathcal{F}(x_j^{[l]},k) + v(x_j^{[l]}, k) \\
    &= \bm{\psi}^T(x_j^{[l]})\bm{\phi}(k) + v(x_j^{[l]}, k) \ \forall\ l\in\{1,\ 2,\ \hdots,\ n_s\}
    \label{eq:indv-agent-single-loc-obs-model}
\end{align}
where $v(x_j^{[l]}, k)$ is some zero-mean additive white Gaussian noise.

Note that we implicitly define the locations as dependent on the timestep, i.e. we have written $x_j^{[l]}(k)$ simply as $x_j^{[l]}$. We can expand this definition to multiple sampling locations in a straightforward manner, giving us:

\begin{equation}
    \mathbf{z}_j(k) = \begin{bmatrix} \bm{\psi}(x_j^{[1]}) & \bm{\psi}(x_j^{[2]}) & \cdots & \bm{\psi}(x_j^{[n_s]})\end{bmatrix}^T\bm{\phi}(k) + \mathbf{v}(k)
    \label{eq:indv-agent-multi-loc-obs-model}
\end{equation}
where we define $\mathbf{v}(k)\sim\mathcal{N}(\mathbf{0}, \mathbf{R}(k))$, with $\mathbf{R}(k)$ as the covariance of the observation noise.

We will define this matrix of Gaussian functions applied at the sampling locations as the observation's process $\mathbf{H}$:

\begin{equation}
    \mathbf{H}(k) = \begin{bmatrix} \bm{\psi}(x_j^{[1]}) & \bm{\psi}(x_j^{[2]}) & \cdots & \bm{\psi}(x_j^{[n_s]})\end{bmatrix}^T
\end{equation}

We now define a centralized version of the linear model for observations; the distributed approach will be discussed in a forthcoming section.

\begin{align}
    \mathbf{z}(k) &= \begin{bmatrix} \mathbf{z}_1^T(k) & \mathbf{z}_2^T(k) & \cdots & \mathbf{z}_{n_r}^T(k) \end{bmatrix}^T \\
    \mathbf{H}(k) &= \begin{bmatrix} \mathbf{H}_1^T(k) & \mathbf{H}_2^T(k) & \cdots & \mathbf{H}_{n_r}^T(k) \end{bmatrix}^T \\
    \mathbf{z}(k) &= \mathbf{H}(k)\bm{\phi}(k) + \mathbf{v}(k)\\
    s.t.\quad \mathbf{H}(k)\bm{\phi}(k) &= \begin{bmatrix} \left(\mathbf{H}_1(k)\bm{\phi}(k)\right)^T & \left(\mathbf{H}_2(k)\bm{\phi}(k)\right)^T & \cdots & \left(\mathbf{H}_{n_r}(k)\bm{\phi}(k)\right)^T\end{bmatrix}^T
\end{align}

We now summarize a linear model for discrete time update of the environment variable $\bm{\phi}(k)$ and the corresponding (centralized) agent observations.

\begin{align}
    \bm{\phi}(k+1) &= \mathbf{A}\bm{\phi}(k) + \mathbf{w}(k) \\
    \mathbf{z}(k) &= \mathbf{H}(k)\bm{\phi}(k) + \mathbf{v}(k)
\end{align}
where $\mathbf{w}(k)\sim\mathcal{N}(\mathbf{0},\mathbf{Q}(k))$, $\mathbf{v}(k)\sim\mathcal{N}(\mathbf{0}_{n_r\times n_s},\mathbf{R}(k))$, and $\mathbf{H}(k)\bm{\phi}(k) =$ \\ $\begin{bmatrix} \left(\mathbf{H}_1(k)\bm{\phi}(k)\right)^T & \left(\mathbf{H}_2(k)\bm{\phi}(k)\right)^T & \cdots & \left(\mathbf{H}_{n_r}(k)\bm{\phi}(k)\right)^T\end{bmatrix}^T$.

% We make $n_s$ measurements from $x_k^{[1]}$ to $x_k^{[n_s]}$, resulting in the observation vector $z^{[1:n_s]}_k=[z_k(x_k^{[1]}),z_k(x_k^{[2]}),\hdots,z_k(x_k^{[n_s]})]^T\in\mathbb{R}^{n_s\times1}$. Going forward, we drop the $k$-subscript indicating time, as we take $n_s$ samples each discrete model time step, giving us $z_k^{[1:n_s]}=[z(x^{[1]}),z(x^{[2]}),\hdots,z(x^{[n_s]})]^T\in\mathbb{R}^{n_s\times1}$. Finally, in the case of $n_r$ agents taking observations, we add in a new subscript $i$ denoting that agent $i$ is taking the observation, i.e. $z_{k,i}=[z_i(x_i^{[1]}),z_i(x_i^{[2]}),\hdots,z_i(x_i^{[n_{s,i}]})]^T\in\mathbb{R}^{n_{s,i}\times1}$ (note that $n_s$ is now changed to also represent the number of samples taken by agent $i$, giving us $n_{s,i}$).

Our goal is to estimate $\bm{\phi}(k)$, and thus $\mathcal{F}(x,k)$, as well as possible; we approach this by attempting to determine an estimate $\hat{\bm{\phi}}(k)$ as close to $\bm{\phi}(k)$ as possible based on some number of agent observations $\mathbf{z}(0:k):=\begin{bmatrix} \mathbf{z}_1(0:k) & \mathbf{z}_2(0:k) & \cdots & \mathbf{z}_{n_r}(0:k) \end{bmatrix}^T$.

\section{Distributed Consensus} \label{sec:distrib-consensus}

\subsection{Linear Filter Updates to the Model: Distributed Kalman Filter} \label{sec:dist-kalman-filt}

\cite{cui2015mutual} shows that we can comfortably realize an observation-based adaptive model update by exploiting the linear structure of the model and implementing a Selective-Basis Kalman Filter. We extend this approach by utilizing the Information variant of the filter, which simultaneously simplifies the distribution process and more explicitly enables sharing of uncertainty values across agents. To start, we quickly recap the application of the standard Kalman Filter towards our centralized formulation. The following are the key values upon which we will filter:

\begin{itemize}
    \item $\hat{\mathbf{x}}_k = \mathbb{E}\left[\mathbf{x}\right]$: The estimate of the state. Traditional literature uses $\hat{\mathbf{x}}_k$ as the standard notation; however, given our model formulation, we will use $\hat{\phi}(k)$.
    \item $\mathbf{P}_{\bm{\phi}}(k)=\mathbb{E}\left[(\bm{\phi}(k)-\hat{\bm{\phi}}(k))(\bm{\phi}(k)-\hat{\bm{\phi}}(k))^T\right]$: The covariance of the stochastic state variable $\bm{\phi}$. We will typically omit the subscript for simplicity.
\end{itemize}

Priors will be notated as $\left(\hat{\bm{\phi}}^-(k), \mathbf{P}^-(k)\right)$.

The filter is a two-step update process, first ``predicting'' a prior estimate on what the filtered state should be, and then ``updating'' the actual value of the state based on observation input. A more detailed description of the math behind the filter is found in \ref{sec:kf-intro-description}. We note that the standard filter requires all observations to be provided prior to update, which is particularly complicated in a distributed setting with intermittent access to neighboring agents. Here, we turn to the Information version of the Kalman Filter, which has an additive update step that allows asynchronous inclusion of novel and new information from neighboring agents. The filter relies on the ``information'' variants of the state variables filtered in the standard filter, defined as in \eqref{eq:information-mat-scribe} and \eqref{eq:information-vec-scribe}, respectively.

\begin{align}
    \mathbf{Y}(k) &= \mathbf{P}_{\bm{\phi}}^{-1}(k) \label{eq:information-mat-scribe} \\
    \mathbf{y}(k) &= \mathbf{P}_{\bm{\phi}}^{-1}(k)\bm{\phi}(k)
    \label{eq:information-vec-scribe}
\end{align}

The priors can be either calculated through the standard Kalman filter and then inverting according to~\eqref{eq:information-mat-scribe} and \eqref{eq:information-vec-scribe}, or through the closed form derivation of the information filter priors seen in \eqref{eq:information-mat-prior-scribe} and \eqref{eq:information-vec-prior-scribe}.

\begin{align}
    \mathbf{M}(k) &= \left(\mathbf{A}^{-1}\right)^T\mathbf{Y}(k-1)\mathbf{A}^{-1} \\
    \mathbf{P}_{\bm{\phi}}(k) &= \mathbf{M}(k) + \mathbf{Q}^{-1}(k) \\
    \mathbf{Y}^{-}(k) &= \mathbf{M}(k) - \mathbf{M}(k)\mathbf{P}_{\bm{\phi}}^{-1}(k)\mathbf{M}(k) \label{eq:information-mat-prior-scribe} \\
    \mathbf{y}^{-}(k) &= \mathbf{Y}^{-}(k)\mathbf{A}\mathbf{Y}(k-1)\mathbf{y}(k-1)
    \label{eq:information-vec-prior-scribe}
\end{align}

With the information filter, we can directly write the update step in the distributed form. Where in our centralized scheme we make the aggregated observation vector $\mathbf{z}(k)$ from individual $\mathbf{z}_j(k)$ observation vectors of each agent, we can instead operate on the individual observation vectors directly for the information update. This gives us \eqref{eq:info-state-innovation-scribe} and \eqref{eq:info-mat-innovation-scribe} to define the innovation provided by observations gathered per agent $j$.

\begin{align}
    \delta\mathbf{i}_j(k) &= \mathbf{H}_j^T(k)\mathbf{R}^{-1}_j(k)\mathbf{z}_j(k) \label{eq:info-state-innovation-scribe} \\
    \delta\mathbf{I}_j(k) &= \mathbf{H}_j^T(k)\mathbf{R}^{-1}_j(k)\mathbf{H}_j(k) \label{eq:info-mat-innovation-scribe}
\end{align}

Given the innovation per agent, we refine our posterior estimate through a simple additive process, as shown in \eqref{eq:info-posterior-state-scribe} and \eqref{eq:info-posterior-cov-scribe}.

\begin{align}
    \mathbf{y}(k) &= \mathbf{y}^-(k) + \sum_{j=1}^{N}\delta\mathbf{i}_j(k) \label{eq:info-posterior-state-scribe} \\
    \mathbf{Y}(k) &= \mathbf{Y}^-(k) + \sum_{j=1}^{N}\delta\mathbf{I}_j(k) \label{eq:info-posterior-cov-scribe}
\end{align}

In the true distributed setting, we make two key changes - observations are no longer aggregated for the linear update, and we average innovations together before summing into the linear update. The former change we have already observed in practice in \eqref{eq:info-state-innovation-scribe} and \eqref{eq:info-mat-innovation-scribe}. The latter is a change we implement to ensure robustness to intermittent communication and noise in the observation process across vehicles. The fundamental implication of averaging is that we are assuming all agents are observing the same environmental process; under this assumption, by averaging we are reducing the innovation to the most common denominator about on the environment behavior that was observed across all agent measurements. In doing so, we smoothen out noise in the observation process, and are additionally able to occasionally miss out on innovations from agents, as they would have only contributed again to the same average value. However, to ensure that the magnitude of the information update is commensurate with the number of agents participant, we scale the innovation addition by the number of agents involved in averaging. This results in the distributed update form seen in \eqref{eq:info-posterior-state-avg-scribe} and \eqref{eq:info-posterior-cov-avg-scribe}.

\begin{align}
    \mathbf{y}(k) &= \mathbf{y}^-(k) + \sum_{j=1}^{N}\frac{1}{N}\sum_{j=1}^{N}\delta\mathbf{i}_j(k) \label{eq:info-posterior-state-avg-scribe} \\
    \mathbf{Y}(k) &= \mathbf{Y}^-(k) + \sum_{j=1}^{N}\frac{1}{N}\sum_{j=1}^{N}\delta\mathbf{I}_j(k) \label{eq:info-posterior-cov-avg-scribe}
\end{align}

Notice that the prior in the distributed information update can be computed locally, with only the information about the system dynamics and the current information estimate. This means that we can express the local prior formation without consideration of other agents or observations, as in Algorithm~\ref{alg:lin-local-update-info-prior}.

\begin{algorithm}
    \caption{Linear Local Update for Information Prior}\label{alg:lin-local-update-info-prior}
    \KwData{Current information estimates $\left[\bm{y}_j,\bm{Y}_j\right](k)$, system descriptors $\left[\bm{A}_j,\ \bm{Q}_j\right](k)$}
    \KwResult{Locally updated priors on information estimates $\left[\bm{y}^-_j,\bm{Y}^-_j\right](k+1)$}
    $\bm{M}_j(k)\gets\left(\bm{A}_j^{-1}(k)\right)^T\bm{Y}_j(k)\bm{A}_j^{-1}(k)$\;
    $\bm{Y}^{-}_j(k)=\bm{M}_j(k)-\bm{M}_j(k)\left(\bm{M}_j(k)+\bm{Q}_j^{-1}(k)\right)^{-1}\bm{M}_j(k)$\;
    $\bm{y}^{-}_j(k)=\bm{Y}^-_j(k)\bm{A}_j(k)\bm{Y}_j(k)\bm{y}_j(k)$\;
\end{algorithm}

\subsubsection{Requirements for Distributed Estimation}

The linear filter update described in the section prior requires two components to successfully work: agents must hold the same prior value $\mathbf{y}^-(k)$ and must jointly average their innovations $\mathbf{i}_j(k)\ \forall\ j$. Distributed averaging of the innovations can be handled through a variety of techniques, as discussed further in \ref{sec:dist-estim-consensus}; SCRIBE utilizes MHMC-based estimators. Agent priors will not diverge so long as agents remain connected to each other throughout operation, guaranteeing that the additive innovations are the same across all agents. However, if communication breaks down, or if not all agents can be reached during an update phase (as is common in asynchronous operation modes), priors can quickly diverge, due to the differences in information available to agents over time. Priors can be quickly brought together via averaging, but as discussed in \ref{sec:cov-intersection}, this runs the risk of overemphasizing certain observations and leading to inconsistent estimates of certainty and information on the environment. We address the fusion of potentially divergent information priors across agents using Hybrid Covariance Intersection.

\subsection{Distributed Estimation Consensus}

The distributed averaging filter builds off of standard centralized averaging filters, whose general form can be expressed as in \eqref{eq:simple-network-averaging-estimator} and is discussed in greater detail in \ref{sec:dist-estim-consensus}. At its core, the approach is to compute a weighted average between the agent's local estimates and that of the agent's neighbors. The weights are often structured such that it assists in speeding up convergence or reducing the likelihood of overconfident estimates.

\begin{equation}
    x_i^{(t+1)}=\sum_{j\in\mathcal{N}_i}w_{ij}x_{j}^{(t)}
    \label{eq:simple-network-averaging-estimator}
\end{equation}

The distributed averaging filter builds off of standard centralized averaging filters, described in greater detail in \ref{sec:dist-estim-consensus}; compared to the centralized scheme, the method presented in \cite{xiao2005scheme} has the additional benefit of not imposing strong requirements on network topology. Instead, it relies entirely on local exchange within a communication graph, between nodes and their immediate neighbors. The nodes achieve consensus by iteratively taking weighted averages of the consensus variables possessed by themselves and their neighbors. The weights are calculated from a Metropolis-Hastings Markov chain (MHMC).

% $\gamma_{kj}^l=\begin{cases} \frac{1}{1+\text{deg}(\mathcal{V}_k)} & \text{if } k\in\mathcal{N}_j \\ 1 - \sum_{k\in\mathcal{N}_j}\gamma_{kj}^l & \text{if } k=j \\ 0 & \text{otherwise} \end{cases}$\;
        % $\overline{\delta\bm{i}}^l_j\gets\sum_{k\in\{j\}\cap\mathcal{N}_j}\gamma_{kj}^l\overline{\delta\bm{i}}_{k}^l$\;
        % $\gamma_{jk}^l=\frac{1}{1+\max\left(\left|\mathcal{N}_j^l\right|,\left|\mathcal{N}_k^l\right|\right)} \forall k\in\mathcal{N}_j^l$\;
        % $\gamma^l_{jj} = \left(1 - \sum_{k \neq j} \gamma^l_{jk}\right) x^l_j$\;
        % $x^{l+1}_i = \gamma^l_{ii} x^l_i + \sum_{j \neq i} \gamma^l_{ij} x^l_j$\;
        % $\overline{\delta\bm{i}}^{l+1}_j \gets \overline{\delta\bm{i}}^l_j + \sum_{k \neq j\in\mathcal{N}_j^l} \gamma^l_{jk} (\overline{\delta\bm{i}}^l_k - \overline{\delta\bm{i}}^l_j)$\;
        % $\overline{\delta\bm{I}}^{l+1}_j = \overline{\delta\bm{I}}^l_j + \sum_{k \neq j\in\mathcal{N}_j^l} \gamma^l_{jk} (\overline{\delta\bm{I}}^l_k - \overline{\delta\bm{I}}^l_j)$

Ultimately, this gives rise to an iterative algorithm, whose update per iteration step can be isolated from previous or future update steps. We can write this update as in Algorithm~\ref{alg:1time-MHMC-linear-filter}.

\begin{algorithm}
    \caption{One-time Update for Distributed MHMC-weighted Linear Filter}\label{alg:1time-MHMC-linear-filter}
    \KwData{Local consensus variable $\bm{x}_i^l$, neighboring consensus variables $\left\{\bm{x}_{j\in\mathcal{N}_i}^l\right\}$}
    \KwResult{MHMC-based update to local variable $\bm{x}_i^{l+1}$}
    $\gamma_{ij}^l\gets\frac{1}{1+\max\left(\left|\mathcal{N}_i^l\right|,\left|\mathcal{N}_j^l\right|\right)} \forall j\in\mathcal{N}_i^l$\;
    $\bm{x}^{l+1}_i \gets \bm{x}^l_i + \sum_{j \neq i\in\mathcal{N}_i^l} \gamma^l_{ij} (\bm{x}^l_j - \bm{x}^l_i)$\;
\end{algorithm}



\subsection{Hybrid Consensus Algorithm} \label{sec:hybrid-consensus}

Compared to a centralized setting, distributed operations more easily result in violations of the Kalman Filter's independence assumption. Consider the example distributed agent network presented in \cite{chen2002estimation}; in this scenario, we have two agents \textit{A} and \textit{B}. Under a distributed filtering scheme, agents each pass the information they have gathered along to their neighbors. Let us then consider the edge-case scenario where \textit{A} receives from \textit{B} information that is not new, but instead originated from \textit{A} itself, due to a quirk of the communication network's topology. Under a Kalman Filtering scheme (for the distributed version, refer to the later section on the distributed Kalman Filter \ref{sec:dist-kalman-filt}), the information matrix of the resulting state estimate $\mathbf{Y}(k)$ will still be increased from the prior, indicating an improvement of knowledge about the state. In reality, however, given that \textit{A} has not considered any "innovation", there is no change in information - this equivalently means our estimate of the uncertainty we have in the system approximation is inconsistent with reality. The occurrences of communication-based violations of the independence assumption is exacerbated by intermittent and unreliable communication conditions. 

To reiterate section \ref{sec:cov-intersection}, the goal of covariance intersection is to provide a consistent fusion of different estimates of a random variable when we lack understanding of the correlation between the estimates. When communication networks break down, causing delays in the consensus process, agents will continue to gather observations. As the information values $\left[\bm{y},\bm{Y}\right](k)$ are updated with these new observations, future observations can no longer be guaranteed as independent of each other (due to a violation of the conditional independence requirement in Markovian chains). With the loss of conditional independence, we can no longer use the vanilla distributed Kalman Filter; furthermore, we have no understanding of what the correlation looks like, giving strong motivation to pursue the distributed method \textit{iterative} covariance intersection.

This method, as discussed in \cite{wang2012distributed}, is a simple extension of the original approach. Here, the optimization process is divided into an iterative procedure; we start by initializing the consensus variables via \eqref{eq:iter-ci-initialization}.

\begin{equation}
    \mathpzc{Y}_j^0\equiv\bm{Y}_j^-(k) \qquad \mathpzc{y}_j^0\equiv\bm{y}_j^-(k)
    \label{eq:iter-ci-initialization}
\end{equation}

After initialization, we solve the optimization subproblem in \eqref{eq:iter-ci-opt-subproblem} each iteration (denoted by superscript $l$). The objective function of the subproblem $\mathcal{J}$ is open for choice, though typically selected as a function of the covariance matrix to provide a measurement of estimate uncertainty, such as $\texttt{trace}(\cdot)$.

\begin{align}
    &\bm{\omega}^*=\arg\min_{\left\{\omega_i\right\}}\mathcal{J}\left(\sum_{i\in\mathcal{N}_j^l}\omega_i\left[\mathpzc{Y}_i^l\right]^{-1}\right) \nonumber \\
    &\text{s.t. }\sum_{i=1}^{\left|\mathcal{N}_j^l\right|}\omega_i=1\ , \qquad\forall i \quad \omega_i \geq 0
    \label{eq:iter-ci-opt-subproblem}
\end{align}

After each iteration update, the local estimate of the variable under consensus is calculated as seen in \eqref{eq:iter-ci-est-update}.

\begin{equation}
    \left[\mathpzc{Y}_j^{l+1},\mathpzc{y}_j^{l+1}\right]=\sum_{i\in\mathcal{N}_j^l}\omega^*_i\left[\mathpzc{Y}_j^l,\mathpzc{y}_j^l\right]
    \label{eq:iter-ci-est-update}
\end{equation}

Given that it is iterative, each update can be written as a procedure, similar to \ref{alg:1time-MHMC-linear-filter}, giving Algorithm~\ref{alg:1time-iter-cov-intersection}.

\begin{algorithm}
    \caption{One-time Update for Distributed Iterative Covariance Intersection}\label{alg:1time-iter-cov-intersection}
    \KwData{Set of neighboring consensus variables $\left\{\mathpzc{Y}_{i\in\mathcal{N}_j}^l,\mathpzc{y}_{i\in\mathcal{N}_j}^l\right\}$}
    \KwResult{Agent-local iterative estimate fusing neighboring information $\left[\mathpzc{Y}_j^{l+1},\mathpzc{y}_j^{l+1}\right]$}
    Solve subproblem \eqref{eq:iter-ci-opt-subproblem} given variables $\left\{\mathpzc{Y}_{i\in\mathcal{N}_j}^l,\mathpzc{y}_{i\in\mathcal{N}_j}^l\right\}$ for $\bm{\omega}^*$\;
    $\left[\mathpzc{Y}_j^{l+1},\mathpzc{y}_j^{l+1}\right]\gets\sum_{i\in\mathcal{N}_j^l}\omega^*_i\left[\mathpzc{Y}_j^l,\mathpzc{y}_j^l\right]$\;
\end{algorithm}

\section{Algorithm Description}

It is helpful for us to restate our purpose when constructing the algorithmic procedure: we wish to specify how multiple autonomous agents, capable of independently gathering measurements of a dynamic environment, can learn a model of the observable environment in such a way that combines the information possessed by all agents to arrive at one shared understanding. We further specify that we wish to do so in a distributed manner, robust to intermittent breaks in the communication network between agents. We divide our goal into three (dependent) sub-tasks, performed by each agent:

\begin{enumerate}
    \item Developing a local model of the environment, using locally gathered observations
    \item Broadcasting locally-learned information and propagating shared information to neighbors in the communication network until it reaches all agents reachable via the network
    \item Update local model parameters such that the resultant model reaches consensus with the local models of all other agents
\end{enumerate}

These three steps are inter-dependent. The information shared in step 2 is determined in step 1, while the local model learned in step 1 will be refined by the output from step 3; the steps required to achieve consensus in step 3 is determined by the state of the communication network in step 2. Of the three steps, we will not explicitly describe the procedure for the second step. However, we will present the two variations of the third step, selected based on the outcome of the second step.

We first develop a procedural approach to environment modeling, our first step, using our understanding from \ref{sec:env-model-formula}, giving us Algorithm~\ref{alg:scal-gaus-field-model}.

\begin{algorithm}
\caption{Distributed Scalar Field Modeling of Environment at Agent $j$}\label{alg:scal-gaus-field-model}
\KwData{Scalar field parameters $\mu_i$, $\sigma_i$, $\tau_i$ for all $i$ in $n_{\phi}$}
\KwResult{Approximation of true environment model $\hat{\mathcal{F}}_j(x,k)$ at agent $j$}
$\psi_i(x) \gets \frac{1}{\tau_i}\exp\left(-\frac{||x-\mu_i||^2_2}{\sigma_i^2}\right)$\;
$k \gets 0$\;
$\bm{\phi}_jj(0) \gets \mathbf{0}_{n_{\phi}}$\;
$\hat{\mathcal{F}}(x,k=0) \gets \bm{\psi}^T(x)\bm{\phi}_j(0)$ \tcc*[r]{Initialize scalar field model}
$\bm{y}_j(0)\gets\bm{0}_{n_\phi}$\;
$\bm{Y}_j(0)\gets\bm{0}_{n_\phi\times n_\phi}$ \tcc*[r]{Initialize information filter values}
\While{Executing Sampling Mission}{
  $\mathbf{z}_j(k) \gets $\textit{GatherSamples}()\;
  $\left[\bm{y}_j,\bm{Y}_j\right](k+1)\gets\text{\textit{DistributedUpdate}}(\bm{\psi},\bm{y}_j(k),\bm{Y}_j(k),\bm{z}_j(k))$\;
  $\bm{\phi}(k+1)\gets\bm{Y}_j^{-1}(k+1)\bm{y}_j(k+1)$\tcc*[r]{Recover stochastic update to $\phi$}
  $\hat{\mathcal{F}}_j(x,k+1) \gets \bm{\psi}^T(x)\bm{\phi}(k+1)$ \tcc*[r]{Update approximate field model}
  $k\gets k+1$\;
}
\end{algorithm}

The procedure is simple, with the bulk of the work contained within the two sub-routines \textit{GatherSamples} and \textit{DistributedUpdate}. The former sub-routine is the algorithm discussed at length in chapter~\ref{sec:dist-task-alloc}. The latter sub-routine refers to our third step in the sub-task description above. In the case where communication has been maintained across agents since the previous update, we are presented with the simpler method of consensus, the standard Information variant of the Kalman Filter, presented in Algorithm~\ref{alg:info-kf-update}. However, in the case that communication has not been maintained, we instead introduce the end result discussed in ~\ref{sec:hybrid-consensus}, shown in Algorithm~\ref{alg:hybrid-consensus-update}. Note that because both approaches utilize the Information variant of the Kalman Filter (or some modification thereof), the scalar field modeling procedure in Algorithm~\ref{alg:scal-gaus-field-model} also maintains the information states for the filter.

\begin{algorithm}
    \caption{Distributed Update for Agent $j$ under stable networks}\label{alg:info-kf-update}
    \KwData{Model basis function $\bm{\psi}$, information values $\left[\bm{y}_j,\bm{Y}_j\right](k)$, and observations $\bm{z}_j$}
    \KwResult{Optimal update to $\left[\bm{y}_j,\bm{Y}_j\right](k+1)$ fusing distributed observations $\bm{z}_{1:n_a}(k)$}
    $\left[\bm{y}_j^-,\bm{Y}_j^-\right](k+1)\gets\text{\textit{LocalPriorUpdate}}\left[\bm{y}_j,\bm{Y}_j\right](k)$\;
    $\bm{H}_j(k)\gets \begin{bmatrix} \psi_1(\bm{z}_j) & \psi_2(\bm{z}_j)& \cdots & \psi_{n_{\phi}}(\bm{z}_j) \end{bmatrix}(k)$ \tcc*[r]{Construct Jacobian}
    $\delta\bm{i}_j\gets\bm{H}_j^T(k)\bm{R}_j^{-1}(k)\bm{z}_j(k)$\;
    $\delta\bm{I}_j\gets\bm{H}_j^T(k)\bm{R}_j^{-1}(k)\bm{H}_j(k)$ \tcc*[r]{Determine locally observed innovation}
    $l \gets 0$\;
    $\left[\overline{\delta\bm{i}}_j^0,\overline{\delta\bm{I}}_j^0\right]\gets\left[\delta\bm{i}_j,\delta\bm{I}_j\right]$\tcc*[r]{Initialize consensus variables}
    \While{Not Converged}{
        \textit{Broadcast}$\left(\left[\overline{\delta\bm{i}}_j^l,\overline{\delta\bm{I}}_j^l\right]\right)$\;
        $\mathcal{M}_j^l\gets$\textit{Collect}$\left(\left\{\overline{\delta\bm{i}}^l_{m\in\mathcal{N}_j},\overline{\delta\bm{I}}^l_{m\in\mathcal{N}_j}\right\}\right)$\;
        $\left[\overline{\delta\bm{i}}_j^{l+1},\overline{\delta\bm{I}}_j^{l+1}\right]\gets$ \textit{MHMC}$\left(\left[\overline{\delta\bm{i}}_j^l,\overline{\delta\bm{I}}_j^l\right], \mathcal{M}_j^l\right)$ \tcc*[r]{MHMC-based local update}
        $l\gets l+1$\;
    }
    $\bm{Y}_j(k+1)\gets\bm{Y}_j^-(k+1)+n_a\overline{\delta\bm{I}}^l_j$\;
    $\bm{y}_j(k+1)\gets\bm{y}_j^-(k+1)+n_a\overline{\delta\bm{i}}^l_j$\tcc*[r]{Update information variables}
\end{algorithm}

The hybrid approach, strongly based on the method introduced in \cite{tamjidi2021unifying}, blends together covariance intersection, distributed Kalman Information Filters, and Gaussian scalar field models, for a reliably distributable model update.

\begin{algorithm}
    \caption{SCRIBE: Distributed Update for Agent $j$ under unstable networks}\label{alg:hybrid-consensus-update}
    \KwData{Model basis function $\bm{\psi}$, information values $\left[\bm{y},\bm{Y}\right](k)$, and observations $\bm{z}_j$}
    \KwResult{Consistent update to $\left[\bm{y}_j,\bm{Y}_j\right](k+1)$ fusing distributed observations $\bm{z}_{1:n_a}$}
    $\left[\bm{y}_j^-,\bm{Y}_j^-\right](k+1)\gets\text{\textit{LocalPriorUpdate}}\left[\bm{y}_j,\bm{Y}_j\right](k)$\;
    $\bm{H}_j(k)\gets \begin{bmatrix} \psi_1(\bm{z}_j) & \psi_2(\bm{z}_j)& \cdots & \psi_{n_{\phi}}(\bm{z}_j) \end{bmatrix}(k)$ \tcc*[r]{Construct Jacobian}
    $\delta\bm{i}_j\gets\bm{H}_j^T(k)\bm{R}_j^{-1}(k)\bm{z}_j(k)$\;
    $\delta\bm{I}_j\gets\bm{H}_j^T(k)\bm{R}_j^{-1}(k)\bm{H}_j(k)$ \tcc*[r]{Determine locally observed innovation}
    $l \gets 0$\;
    $\left[\mathpzc{y}_j^0,\mathpzc{Y}_j^0\right]\gets\left[\bm{y}^-_j,\bm{Y}^-_j\right](k+1)$\;
    $\left[\overline{\delta\bm{i}}_j^0,\overline{\delta\bm{I}}_j^0\right]\gets\left[\delta\bm{i}_j,\delta\bm{I}_j\right]$\tcc*[r]{Initialize consensus variables}
    \While{Not Converged}{
        \textit{Broadcast}$\left(\left[\overline{\delta\bm{i}}_j^l,\overline{\delta\bm{I}}_j^l\right]\right)$\;
        $\mathcal{C}_j^l\gets$\textit{Collect}$\left(\left\{\mathpzc{y}^l_{m\in\mathcal{N}_j},\mathpzc{Y}^l_{m\in\mathcal{N}_j}\right\}\right)$;\quad
        $\mathcal{M}_j^l\gets$\textit{Collect}$\left(\left\{\overline{\delta\bm{i}}^l_{m\in\mathcal{N}_j},\overline{\delta\bm{I}}^l_{m\in\mathcal{N}_j}\right\}\right)$\;
        $\left[\mathpzc{y}_j^{l+1},\mathpzc{Y}_j^{l+1}\right]\gets$ \textit{CI}$\left(\left[\mathpzc{y}_j^{l},\mathpzc{Y}_j^{l}\right], \mathcal{C}_j^l\right)$ \tcc*[r]{CI-based fusion of priors}
        $\left[\overline{\delta\bm{i}}_j^{l+1},\overline{\delta\bm{I}}_j^{l+1}\right]\gets$ \textit{MHMC}$\left(\left[\overline{\delta\bm{i}}_j^l,\overline{\delta\bm{I}}_j^l\right], \mathcal{M}_j^l\right)$ \tcc*[r]{MHMC-based local update}
        $l\gets l+1$\;
    }
    $\bm{Y}_j(k+1)\gets\mathpzc{Y}^l_j+n_a\overline{\delta\bm{I}}^l_j$\;
    $\bm{y}_j(k+1)\gets\mathpzc{y}^l_j+n_a\overline{\delta\bm{i}}^l_j$\;
\end{algorithm}

\section{Results and Remarks}
\subsection{Validation Setup}

We simulate three to five autonomous underwater vehicles (AUVs) exploring and observing a simulated environment. The environment consists of measurable scalar phenomena whose distribution mode varies between being distributed as a Gaussian Scalar Field and being distributed according to an arbitrary spline function with an additive noise process. The agents each take measurements of the environment using imperfect sensors, where the imperfections are modeled as additive zero-mean noise processes. The initial network graph prior to start of operation will be provided, detailing the graph of agents within communicable distance; we assume this is perfectly maintained throughout operation, although the links in the graphs can arbitrarily and randomly change. Communication is represented as a stochastic process, with a scalar parameter describing the rate at which the likelihood of successful communication falls with respect to distance.

% Communication is represented as a stochastic process, with a scalar parameter describing the likelihood of success or failure for \textit{each} round of communication. Dissolved carbon dioxide is simulated by training multiple Gaussian Processes with varying timescales on actual dissolved CO\textsubscript{2} data obtained from the BCO-DMO (Biological \& Chemical Oceanography Data Managament Office) database; the sum of the trained processes is used as an approximation of realistic chemical levels.

Each agent is simulated operating on prefixed lawnmower paths to exhaustively cover the environment; agents operate in non-overlapping regions of the world. Model learning occurs as specified by the SCRIBE algorithm; models are evaluated on prediction accuracy and on distribution distance to the underlying distribution. The former is evaluated by predicting environment states at 100 points equally spaced across the entire environment in a grid pattern, and comparing the model predictions to the output of the underlying environment model. The latter is computed on top of the prior gridpoint predictions, by computing the Frechet distance.

% \noindent\textbf{Experiment Aim}: We will demonstrate that our algorithm enables effective and accurate distributed environment modeling that avoids overconfidence in estimates of uncertainty. We will evaluate the success of our method by comparing the accuracy and confidence of the method against several other approaches, as well as against the ground truth.

% \noindent\textbf{Comparisons}: We consider a variety of approaches for comparison. First is against the ground truth data, to evaluate true accuracy of our modeling technique. To provide upper bounds on estimate accuracy, we will further consider two centralized approaches that assume perfect communication. The first will use a Kalman Filter update on the linear Gaussian scalar field to model the environment, as presented in our work; the second will use a zero-mean RBF (radial basis function) Gaussian Process update for modeling. Finally, to evaluate the effectiveness of the hybrid Kalman Filter-Covariance Intersection technique in distributed modeling, we will also compare against two distributed approaches that do not assume perfect communication. One approach will only use a Kalman Filter to update the distributed Gaussian scalar field model, while the other will only use Covariance Intersection updates for the same.

% \noindent\textbf{Performance at varying levels of communication}: We will set out to demonstrate the effectiveness of our technique at \textit{any level of communicability} by repeating the modeling and comparison experiments for multiple values for the communicability parameter. We will range from perfect communication to no communication.

% \noindent\textbf{Other assumptions about the agent and environment}: We assume that we can safely ignore the effects of disturbances on agent operation beyond communication; examples include sensor failure and drift, motor and position drift, and persistent environmental occlusion of sensors. We posit that this is a responsibility of the agent to address outside of the modeling process.

\subsection{Results}

Figures~\ref{fig:scribe-res-basic-perf}, \ref{fig:scribe-no-comms-v-comms-accuracy}, and \ref{fig:scribe-accuracy-over-time}, show the results of the simulated trials. Overall, the model learning technique provided by SCRIBE shows high accuracy, providing up to 95\% accuracy while under perfect communication and up to 90\% accuracy even under intermittent commmunication. Individual agent models remained consistent with each other even under intermittent communication. These takeaways indicate that the distributed model learning technique does indeed function successfully and robustly to communication flaws. As the learning technique is ultimately a linear model formed by composing Gaussians, we find the increase in model prediction error at phenomena measurement peaks to be consistent with intuition. Furthermore, we note that the learned model still sufficiently captures the phenomena distribution, i.e. the error in learned model predictions is not simply due to the phenomena deviating from a constant mean value.

\begin{figure}[htb]
    \centering
    \includegraphics[width=0.75\linewidth]{images/chapter3/init_results_scribe.jpg}
    \caption{Simulation results when the agents know \textit{a priori} the locations of sites with high-intensity phenomena, operating under perfect communication.}
    \label{fig:scribe-res-basic-perf}
\end{figure}

\begin{figure}[htb]
    \centering
    \begin{subfigure}[t]{0.45\textwidth}
        \centering
        \includegraphics[width=\columnwidth]{images/chapter3/3a_121w_fullGT_nonec_lowo_10s_lawnmower_err_map.png}
        \caption{Accuracy of SCRIBE's distributed model after ten rounds of iteration, with no communication between agents.}
        \label{fig:scribe-no-comms-10}
    \end{subfigure}
    \begin{subfigure}[t]{0.45\textwidth}
        \centering
        \includegraphics[width=\columnwidth]{images/chapter3/3a_121w_fullGT_distc_lowo_10s_lawnmower_err_map.png}
        \caption{Accuracy of SCRIBE's distributed model after ten rounds of iteration, with communication between agents each round.}
        \label{fig:scribe-w-comms-10}
    \end{subfigure}
    \caption{Comparison between SCRIBE's model's performance under different communication conditions. Both show strong accuracy, while SCRIBE's model is able to reduce inaccuracy across the board noticeably more when communication is allowed.}
    \label{fig:scribe-no-comms-v-comms-accuracy}
\end{figure}

\begin{figure}[htb]
    \centering
    \begin{subfigure}[t]{0.45\textwidth}
        \centering
        \includegraphics[width=\columnwidth]{images/chapter3/3a_121w_fullGT_distc_lowo_10s_lawnmower_err_map.png}
        \caption{SCRIBE's distributed model accuracy after ten rounds of iteration, with communication between agents.}
        \label{fig:scribe-w-comms-10-2}
    \end{subfigure}
    \begin{subfigure}[t]{0.45\textwidth}
        \centering
        \includegraphics[width=\columnwidth]{images/chapter3/3a_121w_fullGT_distc_lowo_20s_lawnmower_err_map.png}
        \caption{SCRIBE's distributed model after twenty rounds of iteration, with communication between agents each round.}
        \label{fig:scribe-w-comms-20}
    \end{subfigure}
    \caption{Comparison between SCRIBE's model's performance over time. Both show strong accuracy, while SCRIBE's model improves as more samples are gathered over time.}
    \label{fig:scribe-accuracy-over-time}
\end{figure}

We further see that when agents are provided a general set of Gaussian scalar field bases for model learning, as opposed to knowing \textit{a priori} where phenomena are expected to occur, the accuracy remains high. This is further supported under simulation modes where the underlying phenomena are modeled not via Gaussian scalar fields but by spline functions with added noise, where the learned model's predictive accuracy continues to be high. Finally, we note that the distance between the learned model's predictive distribution and the true phenomena distribution decreases across iterations, indicating stable learning performance.

\subsection{Concluding Takeaways}

In summary, the distributed model learning technique provided through SCRIBE shows strong potential for enabling effective and efficient online environment model learning. The learning technique is computationally light, and  robust to intermittent communication, maintaining learned model accuracy and consistency even in the face of occassional disconnections between agents. SCRIBE's approach achieves learned model with predictive accuracy as high as 90\% even under imperfect communication conditions, and is usable by agents under asynchronous operational modes.

It is important to note that as we utilized a linear model formulation with fixed hyper-parameters (Gaussian scalar field parameters) in SCRIBE, the learned model is less flexible in capturing world phenomena compared to other approaches. However, the linear formulation further enables low-cost, highly compressed communication for model learning. Furthermore, this linear model is comparable to a fixed-size Gaussian Process with a Radial Basis Function, and thus the hyper-parameters can be intelligently determined via a Gaussian Process model learned offline with prior information. Finally, we believe this mechanism for selecting the Gaussian scalar field parameters can be improved with the inclusion of \textit{E}-dimensional linear approximations of Gaussian Processes \cite{pillonetto2018distributed}, as specifically formulated for multi-agent dynamic sensor networks in \cite{jang2021fully, kim2023distributed}. These improvements are left as future directions of improving the SCRIBE formulation presented in this chapter.